\documentclass[a4paper, noanswers, addpoints, 12pt]{exam}
\usepackage[T1]{fontenc}
\usepackage[utf8]{inputenc}
\usepackage{lmodern}
\usepackage[french]{babel}
\usepackage{xcolor}
\usepackage{listings}
\usepackage{graphicx}
\usepackage[most]{tcolorbox}

\definecolor{codegreen}{rgb}{0,0.6,0}
\definecolor{codegray}{rgb}{0.5,0.5,0.5}
\definecolor{codepurple}{rgb}{0.58,0,0.82}
\definecolor{backcolour}{rgb}{0.95,0.95,0.92}
 
\lstdefinestyle{mystyle}{
    backgroundcolor=\color{backcolour},   
    commentstyle=\color{codegreen},
    keywordstyle=\color{magenta},
    numberstyle=\tiny\color{codegray},
    stringstyle=\color{codepurple},
    basicstyle=\footnotesize,
    breakatwhitespace=false,         
    breaklines=true,                 
    captionpos=b,                    
    keepspaces=true,                 
    numbers=left,                    
    numbersep=5pt,                  
    showspaces=false,                
    showstringspaces=false,
    showtabs=false,                  
    tabsize=2
}
\lstset{style=mystyle, language=Python}

\title{CC1 Prem NSI - Conditions et boucles}
\author{Babash - NSI Tani Malandi et Johanne Bakalara}
\date{Octobre 2024, Durée: 1h}


\usepackage[a4paper,top=2.3cm,bottom=2.8cm,left=2cm,right=2cm,headsep=0.6cm,footskip=1.2cm]{geometry}
\lstset{extendedchars=true,literate={é}{{\'e}}1}
\begin{document}
\runningfooter{}{\footnotesize Babash - NSI Tani Malandi et Johanne Bakalara \textperiodcentered{} CC BY-SA 4.0 \textperiodcentered{} page \thepage}{}
\firstpagefooter{}{\footnotesize Babash - NSI Tani Malandi et Johanne Bakalara \textperiodcentered{} CC BY-SA 4.0 \textperiodcentered{} page \thepage}{}

\maketitle
\thispagestyle{headandfoot}

\pagestyle{headandfoot}
\thispagestyle{headandfoot}
\firstpageheader{\makebox[\textwidth]{Nom : \dotfill\ Pr\'enom : \dotfill\ Classe : \makebox[3cm]{\dotfill}}}{}{}


\section*{Consignes:}
\begin{center} 
\fbox{\parbox{\linewidth}{
\begin{center}
Le devoir est à rédiger sur une feuille libre.
% La clarté de la rédaction et les explications font partie intégrante de la notation. Une réponse non rédigée ou sans explication sera considérée fausse.\\
Prenez le temps d'encadrer les réponses.
\end{center}
}}
\end{center}

\vsword{Note}
\hsword{Note}

\section*{Barème :}
\begin{center}
  \par\bigskip
  \gradetable[h]
\end{center}
\begin{questions}

\section{Question de cours : }

\question[2]  :
\begin{parts}
    \part Quels sont les 4 types primitifs en Python ?
    \part Quelle fonction permet d'afficher une information d'un utilisateur ?
    \part Quelle fonction permet de récupérer la saisie d'un utilisateur en Python ?
    \part Que signifie le terme indentation ?
\end{parts}

\question[2] Donner les plages de valeurs correspondant aux appels de range suivants :
\begin{itemize}
    \item range(10)
    \item range(5, 7)
    \item range(7, 17, 2)
    \item range(98, 0, -10)
\end{itemize}

\question[2] Quels opérateurs Python permettent de réaliser les tests suivants :
\begin{itemize}
    \item différent de :
    \item strictement supérieur à :
    \item inférieur ou égal à :
    \item supérieur ou égal à :
\end{itemize}
 
\section{Variables et conditions : }

\begin{lstlisting}[language=Python]
    import random
    # random.randint(1, 10) tire un nombre aléatoire entre 1 et 10.
    carte1 = random.randint(1, 10)
    carte2 = random.randint(1, 10)
    carte3 = random.randint(1, 10)
    if carte1 == carte2 == carte3 == 1:
        score = 300
    elif carte1 == carte2 == carte3:
        score = 30
    elif carte1 == carte2 or carte2 == carte3 or carte1 == carte3:
        score = 20
    else:
        score = carte1 + carte2 + carte3
    if carte1 == 10 or carte2 == 10 or carte3 == 10:
        score = score + 10
\end{lstlisting}
\question[1] Combien y a-t-il de variables dans ce code, quelles sont leur type ?
\question[3] Donner le tableau d'avancement des variables pour carte1=1, carte2=1, carte3=3.
\question[2] Quel est le score pour :
\begin{itemize}
    \item carte1=1, carte2=1, carte3=1
    \item carte1=10, carte2=10, carte3=3
    \item carte1=10, carte2=9, carte3=9
    \item carte1=1, carte2=9, carte3=2
\end{itemize}

\section{Boucles : }

\question[1] Combien y a-t-il de variables dans ce code, quelles sont leur type ?
\question[4] Faire le tableau d'avancement des variables du code ci-dessous :
\begin{lstlisting}[language=Python]
n = 3
for i in range(n):
   a = i
   while a > 1:
      a = a/2
\end{lstlisting}

\section{Rédaction de code : }
\question[3] Écrire un code permettant de calculer la somme des multiples entiers de 3 entre 1 et 10 000. On pourra utiliser l'opérateur (\%).


\end{questions}

\end{document}

